% Options for packages loaded elsewhere
\PassOptionsToPackage{unicode}{hyperref}
\PassOptionsToPackage{hyphens}{url}
%
\documentclass[
  11pt,
]{article}
\usepackage{amsmath,amssymb}
\usepackage{iftex}
\ifPDFTeX
  \usepackage[T1]{fontenc}
  \usepackage[utf8]{inputenc}
  \usepackage{textcomp} % provide euro and other symbols
\else % if luatex or xetex
  \usepackage{unicode-math} % this also loads fontspec
  \defaultfontfeatures{Scale=MatchLowercase}
  \defaultfontfeatures[\rmfamily]{Ligatures=TeX,Scale=1}
\fi
\usepackage{lmodern}
\ifPDFTeX\else
  % xetex/luatex font selection
\fi
% Use upquote if available, for straight quotes in verbatim environments
\IfFileExists{upquote.sty}{\usepackage{upquote}}{}
\IfFileExists{microtype.sty}{% use microtype if available
  \usepackage[]{microtype}
  \UseMicrotypeSet[protrusion]{basicmath} % disable protrusion for tt fonts
}{}
\makeatletter
\@ifundefined{KOMAClassName}{% if non-KOMA class
  \IfFileExists{parskip.sty}{%
    \usepackage{parskip}
  }{% else
    \setlength{\parindent}{0pt}
    \setlength{\parskip}{6pt plus 2pt minus 1pt}}
}{% if KOMA class
  \KOMAoptions{parskip=half}}
\makeatother
\usepackage{xcolor}
\usepackage[margin=24mm]{geometry}
\setlength{\emergencystretch}{3em} % prevent overfull lines
\providecommand{\tightlist}{%
  \setlength{\itemsep}{0pt}\setlength{\parskip}{0pt}}
\setcounter{secnumdepth}{-\maxdimen} % remove section numbering
\ifLuaTeX
  \usepackage{selnolig}  % disable illegal ligatures
\fi
\usepackage{bookmark}
\IfFileExists{xurl.sty}{\usepackage{xurl}}{} % add URL line breaks if available
\urlstyle{same}
\hypersetup{
  pdftitle={Two-moment control of the original arithmetic relation blocks},
  pdfauthor={Owner-directed SplitZero analytic continuation},
  hidelinks,
  pdfcreator={LaTeX via pandoc}}

\title{Two-moment control of the original arithmetic relation blocks}
\usepackage{etoolbox}
\makeatletter
\providecommand{\subtitle}[1]{% add subtitle to \maketitle
  \apptocmd{\@title}{\par {\large #1 \par}}{}{}
}
\makeatother
\subtitle{An upper determinant certificate without a separately supplied
restriction gap}
\author{Owner-directed SplitZero analytic continuation}
\date{13 September 2026}

\begin{document}
\maketitle

{
\setcounter{tocdepth}{3}
\tableofcontents
}
\subsection{1. Result and the current inherited
theorem}\label{result-and-the-current-inherited-theorem}

This note continues the canonical endpoint-restriction construction. It
does not choose a different positive metric or replace the arithmetic
source by a comparison measure. Its new application is an upper bound
for each original relation-block volume from two source-derived traces,
together with a complete positive remainder and a source interpretation
of the second trace as exterior area.

The general sharp determinant optimization with prescribed dimension,
trace, and squared trace is classical; Grone, Johnson, Marques de Sá,
and Wolkowicz (1984) study precisely that optimization. The proof below
is included in full, using a scalar Hermite interpolation identity. No
priority claim is made for the general inequality. The work here is its
composition with the actual arithmetic representative, internal
quotient, inverse restriction, exterior source maps, and
consecutive-window theorem.

For one original block, the output is

\[
 \boxed{\frac{V_i}{V_j}\le(1+a_{i,j})(1+b_{i,j})^{q-1}.}                 \tag{1}
\]

Both \(a_{i,j}\) and \(b_{i,j}\) are calculated from the first two
traces of a specified endomorphism. There is no assumption that the
arithmetic action is normal, no assumption that a restriction eigenspace
is invariant under it, and no independently supplied lower gap for the
restriction.

\subsubsection{Integration of the two new
reports}\label{integration-of-the-two-new-reports}

The attachment reporting PR \#25 supplies the phase-retaining
consecutive product, the separate first-admissible-degree step, a
two-sided written arithmetic norm envelope, and its comparable-window
consequence. Its resulting off-line-quartet lower budgets are

\[
 \log(V_q/V_{2q})\ge2q\log(D_hk),\qquad
 \log(V_{q-1}/V_{2q-1})\ge2q\log(D_hk),                         \tag{2}
\]

for a fixed positive \(D_h\), hence the four-volume lower threshold is
\textbf{four}, not merely the earlier sufficient lower threshold two.
This is inherited from the parallel contributions, not claimed as a
discovery of this note. The source explicitly distinguishes its checked
finite implications from the written analytic estimates and the
asymptotic conclusion.

Its regularity clarification is retained: the actual arithmetic density
is in \(L^1\cap L^\infty\), and

\[
 \|(\tau_aw_h-w_h)*w_h\|_\infty
 \le\|\tau_aw_h-w_h\|_1\|w_h\|_\infty\longrightarrow0.
\]

Here \(\tau_a\) denotes translation of a function, not the split scalar
\(\tau\); its definition is \((\tau_af)(t)=f(t-a)\). This supplies
continuity of the convolution used in the source's lower-envelope
argument.

The second attachment is a working transcript reporting further Gamma,
coefficient-error, endpoint, and reflected-algebra work. Its reported
moment certificates and unpublished upper-bound details are not
independently reproduced here or silently made premises. In particular,
this note does not duplicate its claimed fixed-budget extremizer or
claim to have executed its local Windows commands.

GitHub main was read at
\texttt{952ef9fee1e1419b6858d920354de8fa99430b7d}. PR \#25 was read at
its reported verified head
\texttt{66fac5e7885a40cd4873e74b754907320f05b067}, including its full
\texttt{RESEARCH\_NOTE.md}. Its CI status is a source report; this
session did not execute Lean.

\subsection{2. Original arithmetic source and
support}\label{original-arithmetic-source-and-support}

Retain the original scalar object and amplitude map

\[
 G(R)=\{\tau\}\sqcup\{r^\bullet:r\in R\},\qquad
 e_R=0_R^\bullet,\qquad
 p_R(\tau)=0_R,\quad p_R(r^\bullet)=r.
\]

The original coefficient square remains

\[
 \begin{array}{ccc}
 G(\mathbb Z)&\xrightarrow{G(j)}&G(\mathbb C)\\
 p_{\mathbb Z}\downarrow&&\downarrow p_{\mathbb C}\\
 \mathbb Z&\xrightarrow{j}&\mathbb C.
 \end{array}                                                     \tag{3}
\]

Its arithmetic target \(\mathbb Z\) is infinite. The geometric base
remains \(\mathfrak b_\tau\). The square does not replace the internal
support fibres by the ordinary amplitude observation.

Keep the original theta complex and generator

\[
 C_+=[V\xrightarrow{\Theta}\mathscr B],\qquad
 Q=\mathscr B/\Theta V,\qquad D=-x\partial_x,\qquad g=2\xi.
\]

For a nonempty full packet \(h\) of actual zeros and tensor degree
\(k\), retain the cyclic module and its action

\[
 E=\mathbb C[S]/(\chi_{h,k}),\quad q=\deg\chi_{h,k},\quad A=M_S,
 \quad S=s_1+\cdots+s_k.
\]

The source and arithmetic inclusion are the constructed maps

\[
 \mathcal VP=P(D_1+\cdots+D_k)(F_h^{\otimes k}),\qquad
 \mathcal MF_h=g/h,
\]

\[
 \eta[P]=\upsilon_h^{\otimes k}P(A_k)1,\qquad
 \upsilon_h=j_h(g/h),                                             \tag{4}
\]

\[
 J^{(k)}\mathcal V=\eta\pi_\chi,\qquad
 q^{(k)}\mathcal V=\sigma_h^{\otimes k}\eta\pi_\chi.
\]

The original norm is the integral against \(m_{h,k}=w_h^{*k}\) on
\(S=k/2+iu\), where

\[
 w_h(t)=\frac{|(g/h)(1/2+it)|^2}{2\pi},\qquad
 \int m_{h,k}=\mu_h^k.
\]

No mass is assigned a new value. All nilpotent jet orders and the unit
in (4) remain in the arithmetic maps.

For \(h=1\) the finite packet is the zero coefficient module. Its
determinant is the empty determinant one, its log-volume is zero, and
\(G(0)\) still contains the two split elements. The analytic theta seed
remains nonzero. Positive-dimensional formulas below use \(q\ge1\); the
divisions by \(q-1\) use \(q\ge2\).

\subsubsection{Fixed polynomial
coordinates}\label{fixed-polynomial-coordinates}

Let \(p_0,p_1,\ldots\) be the original monic orthogonal polynomials and
\(\omega_n>0\) their actual squared norms. Keep \(\omega_0=\mu_h^k\). In
fixed remainder coordinates on \(E\), set

\[
 b_n=[p_n]_\chi,\quad B_N=[b_0,\ldots,b_N],\quad
 O_N=\operatorname{diag}(\omega_0,\ldots,\omega_N).
\]

For every admitted \(N\ge q-1\),

\[
 K_N=B_NO_N^{-1}B_N^*,\quad G_N=K_N^{-1},\quad
 R_N=O_N^{-1}B_N^*G_N,\quad V_N=\det G_N.                         \tag{5}
\]

Thus \(B_NR_N=I\) and \(R_N^*O_NR_N=G_N\). The actual test-function
representative is obtained by the specified polynomial-column
realization \(\mathcal V\), not by treating a coefficient vector as a
function without that map.

\subsection{3. The positive cost operator comes from the original
restriction}\label{the-positive-cost-operator-comes-from-the-original-restriction}

Fix \(q-1\le i\le j\). Set \(r=j-i\) and retain the new columns and
their norms

\[
 F=F_{i,j}=[b_{i+1},\ldots,b_j]:\mathbb C^r\to E,\qquad
 \Omega=\operatorname{diag}(\omega_{i+1},\ldots,\omega_j),
\]

\[
 D_{i,j}=K_j-K_i=F\Omega^{-1}F^*.                                \tag{6}
\]

The Hilbert spaces \(E_i=(E,G_i)\) and \(E_j=(E,G_j)\) are connected by
the actual forward transport \(I_{i,j}:E_i\to E_j\), identity on the
fixed coefficient space. Its adjoint is

\[
 T_{i,j}=I_{i,j}^{\dagger}=K_iG_j:E_j\to E_i.
\]

To form an endomorphism of one specified Hilbert object, compose first:

\[
 Q_{i,j}=T_{i,j}I_{i,j}:E_i\to E_i.
\]

The coefficient matrix of \(Q_{i,j}\) is \(K_iG_j\). Define

\[
 \boxed{Z_{i,j}=Q_{i,j}^{-1}-I_E=(K_j-K_i)G_i=D_{i,j}G_i.}       \tag{7}
\]

It is positive self-adjoint on \(E_i\), because

\[
 G_iZ_{i,j}=G_iD_{i,j}G_i\succeq0.
\]

Its eigenvalues \(\lambda_a\) are nonnegative. The restriction
eigenvalues in the preceding note are precisely

\[
 g_a=\frac1{1+\lambda_a},\qquad
 \frac{V_i}{V_j}=\det(I_E+Z_{i,j}),\qquad
 \log(V_i/V_j)=\sum_{a=1}^{q}\log(1+\lambda_a).                 \tag{8}
\]

No inverse-square-root transformation changes the canonical arithmetic
metric. Equation (7) is a typed inverse of the already-constructed
restriction endomorphism.

\subsubsection{Source and action maps remain
attached}\label{source-and-action-maps-remain-attached}

The source map \(F:(\mathbb C^r,\Omega)\to(E,G_i)\) has adjoint

\[
 F^{\dagger}=\Omega^{-1}F^*G_i,
 \qquad Z_{i,j}=FF^{\dagger}.                                   \tag{9}
\]

Let \(\mathsf H_N=G_N^{-1}(A^*G_N+G_NA-kG_N)\) be the inherited control
endomorphism. Put
\(\widehat{\mathsf H}_j=I_{i,j}^{-1}\mathsf H_jI_{i,j}\) on \(E_i\). The
previous action defect and inversion give

\[
 \boxed{[A,Z_{i,j}]=
 \widehat{\mathsf H}_j(I+Z_{i,j})-(I+Z_{i,j})\mathsf H_i.}        \tag{10}
\]

Indeed \([A,Q]=\mathsf H_iQ-Q\widehat{\mathsf H}_j\) and
\([A,Q^{-1}]=-Q^{-1}[A,Q]Q^{-1}\). This is the actual relation between
the cost operator and the arithmetic generator. A cost eigenspace is not
presumed invariant for \(A\).

\subsection{4. The two input traces are norms of actual source
maps}\label{the-two-input-traces-are-norms-of-actual-source-maps}

Write

\[
 t_1=\operatorname{Tr}Z_{i,j},\qquad
 t_2=\operatorname{Tr}(Z_{i,j}^2),\qquad
 \mathcal A_2=\frac{t_1^2-t_2}{2}.                                      \tag{11}
\]

The scalar \(\mathcal A_2\) is the second exterior coefficient of this
cost operator. Its split observation is
\(\mathcal A_2^\bullet\in G(\mathbb R)\); a vanishing value maps to the
represented scalar \(e\).

From (6) and cyclicity of the finite trace,

\[
 \boxed{t_1=\sum_{n=i+1}^{j}\frac{b_n^*G_ib_n}{\omega_n},}
\]

\[
 \boxed{t_2=\sum_{n,m=i+1}^{j}
 \frac{|b_n^*G_ib_m|^2}{\omega_n\omega_m},}                      \tag{12}
\]

\[
 \boxed{\mathcal A_2=\sum_{i<n<m\le j}
 \frac{(b_n^*G_ib_n)(b_m^*G_ib_m)-|b_n^*G_ib_m|^2}
 {\omega_n\omega_m}.}                                          \tag{13}
\]

The off-diagonal pairings in (12)--(13) are essential.

Let \(\Phi=R_iF\) with its actual degree-\(i\) source norm. Then
\(t_1=\|\Phi\|_{\mathrm{HS}}^2\). Equip unscaled exterior bases with
determinant Grams. The precise second identity is

\[
 \boxed{\mathcal A_2=\|\bigwedge^2\Phi\|_{\mathrm{HS}}^2.}               \tag{14}
\]

Proof: in the coordinate wedge indexed by \(n<m\), the squared source
norm is \(\omega_n\omega_m\). Its image squared norm is the two-by-two
Gram determinant in (13). Summing proves (14). An alternating-tensor
realization instead has its literal factor \(2!\) in each wedge Gram;
this statement uses the specified exterior quotient, so no factorial is
dropped.

Let \(T_{\mathrm{new}}\) denote the actual newly admitted orthogonal
polynomial columns in the degree-\(j\) source. The source relation is

\[
 \mathfrak b_{i,j}=T_{\mathrm{new}}-L_{i,j}R_iF,
 \qquad
 B_j\mathfrak b_{i,j}=0,
\]

\[
 \boxed{\mathfrak b_{i,j}^*O_j\mathfrak b_{i,j}
 =\Omega+F^*G_iF.}                                             \tag{15}
\]

Thus (12)--(14) measure the old representative components needed to
cancel the new columns' jets. The relation itself maps to the receiving
supported zero through the original coequalizer. Neither component is
declared absent.

An admitted pair with dependent images has zero exterior area in (13), a
represented zero. The pair and its original support label remain
available. A pair outside the degree-indexing object is absent. No
statement that these two cases have the same upstream datum is used.

\subsection{5. A sharp two-trace upper bound, with a complete
proof}\label{a-sharp-two-trace-upper-bound-with-a-complete-proof}

Suppose \(q\ge2\). Define

\[
 d=\sqrt{\frac{qt_2-t_1^2}{q-1}},\qquad
 a=\frac{t_1+(q-1)d}{q},\qquad
 b=\frac{t_1-d}{q}.                                            \tag{16}
\]

Both discriminants and signs come from the actual positive spectrum:

\[
 \frac{t_1^2}{q}\le t_2\le t_1^2,
 \qquad 0\le b\le a,
\]

\[
 a+(q-1)b=t_1,\qquad a^2+(q-1)b^2=t_2.                         \tag{17}
\]

\subsubsection{The spectral cap is proved, not
supplied}\label{the-spectral-cap-is-proved-not-supplied}

For any eigenvalue \(x\) of \(Z\), Cauchy--Schwarz on the other \(q-1\)
eigenvalues gives

\[
 t_2\ge x^2+\frac{(t_1-x)^2}{q-1}.
\]

Solving this quadratic yields \(x\le a\). Hence the same two traces also
produce the actual restriction-gap bound

\[
 \boxed{Q_{i,j}\succeq\frac1{1+a}I_E
 \quad\text{in the original metric }G_i.}                      \tag{18}
\]

This is a finite derived bound, not a uniform estimate in \(k\).

\subsubsection{The optimal determinant bound from these two
traces}\label{the-optimal-determinant-bound-from-these-two-traces}

The result is

\[
 \boxed{
 \log\det(I+Z)\le U_2(Z):=\log(1+a)+(q-1)\log(1+b).
 }                                                             \tag{19}
\]

For \(a>b\), define the quadratic

\[
 P_{a,b}(x)=\log(1+b)+\frac{x-b}{1+b}+c_{a,b}(x-b)^2,
\]

\[
 c_{a,b}=
 \frac{\log(1+a)-\log(1+b)-(a-b)/(1+b)}{(a-b)^2}<0.              \tag{20}
\]

It matches \(\log(1+x)\) and its derivative at \(b\), and the value at
\(a\). For \(0\le x\le a\) the exact remainder is

\[
 \boxed{
 P_{a,b}(x)-\log(1+x)
 =(a-x)(x-b)^2
 \int_1^\infty\frac{ds}{(s+x)(s+b)^2(s+a)}\ge0.
 }                                                             \tag{21}
\]

Proof of (21): use \(\log(1+x)=\int_1^\infty(1/s-1/(s+x))\,ds\) and the
exact third divided difference of \(1/(s+x)\). Interpolation at the
repeated point \(b\) and at \(a\) gives the displayed rational
denominator. The integral converges as \(s^{-4}\); all factors have
their shown positive signs. This also directly proves the Hermite
remainder without an unspecified intermediate point.

Sum (21) over the actual eigenvalues. Since a quadratic trace depends
only on \(q,t_1,t_2\), equations (17) give

\[
 \operatorname{Tr}P_{a,b}(Z)=P_{a,b}(a)+(q-1)P_{a,b}(b)=U_2(Z).
\]

This proves (19) and (1).

If \(a=b\), equality in \(qt_2\ge t_1^2\) forces every eigenvalue to
equal \(b\), so (19) is equality without the quotient in (20). If
\(q=1\), the exact answer is \(\log(1+t_1)\). If \(t_1=0\), positivity
gives \(Z=0\) and the volume is zero. If \(q=0\), use the empty
determinant, without forming (16).

For \(a>b\), the integrand in (21) is strictly positive except at
\(x=a,b\). Equality in (19) therefore holds exactly when the spectrum
consists of one \(a\) and \(q-1\) copies of \(b\). These values have the
original two traces by (17), so the bound is best possible given only
those traces and the dimension. This does not assert that every such
extremizer arises from an arithmetic packet.

The general optimization is classical. The source application uses
precisely \(Z=(K_j-K_i)G_i\), not an independently chosen Hermitian
matrix.

\subsection{6. Retain the complete remainder; a third trace improves the
bound}\label{retain-the-complete-remainder-a-third-trace-improves-the-bound}

Let

\[
 C_3=\operatorname{Tr}\bigl((aI-Z)(Z-bI)^2\bigr)
 =-t_3+(a+2b)t_2-(2ab+b^2)t_1+qa b^2,
\]

where \(t_3=\operatorname{Tr}Z^3\). Positivity on the spectrum proves
\(C_3\ge0\).

The full loss in the upper bound is

\[
 \boxed{
 U_2(Z)-\log\det(I+Z)
 =\int_1^\infty
 \frac{\operatorname{Tr}((aI-Z)(Z-bI)^2(sI+Z)^{-1})}
 {(s+b)^2(s+a)}\,ds.
 }                                                             \tag{22}
\]

All four operator factors commute because they are functions of this
same \(G_i\)-self-adjoint \(Z\). The expression is positive in that same
metric. No claim about commutation with \(A\) replaces (10).

Since \(0\le x,b\le a\), the integrand denominator in (21) is at most
\((s+a)^4\). Therefore

\[
 \boxed{\log\det(I+Z)
 \le U_2(Z)-\frac{C_3}{3(1+a)^3}.}                              \tag{23}
\]

This retains a concrete part of the higher-spectrum correction. It
requires one further trace of the same source cost, not a new spectral
gap or arithmetic realization.

Equation (19) is always at least as sharp as the trace-only bound

\[
 U_2(Z)\le q\log(1+t_1/q),                                     \tag{24}
\]

by scalar concavity and (17). Equality in (24), apart from a
zero-dimensional case, requires \(a=b\). It is not asserted to dominate
every higher-moment certificate from the previous note; those can be
retained and their certified upper bounds compared.

\subsection{7. The simplest bound exposes independent pair
growth}\label{the-simplest-bound-exposes-independent-pair-growth}

For \(q\ge2\) and \(t_1>0\), (17) implies

\[
 \mathcal A_2=(q-1)b\left(t_1-\frac q2b\right).
\]

Because \(0\le b\le t_1/q\),

\[
 \frac{\mathcal A_2}{(q-1)t_1}\le b
 \le\frac{2\mathcal A_2}{(q-1)t_1},\qquad a\le t_1.
\]

Consequently there is a radical-free finite upper bound

\[
 \boxed{
 \frac{V_i}{V_j}
 \le(1+t_1)
 \left(1+\frac{2\mathcal A_2}{(q-1)t_1}\right)^{q-1}.
 }                                                             \tag{25}
\]

The quantity \(\mathcal A_2\) is exactly the sum of the original
two-column exterior areas (13), including every cross-pairing. The first
trace alone cannot distinguish nearly parallel old representative
corrections from many independent ones; (13) does make that distinction
quantitatively.

If \(\mathcal A_2=0\) and \(t_1>0\), the positive operator has rank one,
and the sharper formula (19) is the exact value \(\log(1+t_1)\). This is
a statement about the rank of this specified comparison map, not a claim
about uniqueness of the original scalar support or about the number of
higher support fibres.

\subsection{8. Compose with the new consecutive-window
theorem}\label{compose-with-the-new-consecutive-window-theorem}

The inherited phase-retaining product implies for every original
admitted window \(i,\ldots,j-1\), \(r=j-i\), including \(i=q-1\) through
its separately proved first-degree step,

\[
 \min_{i\le N<j}\epsilon_N
 \le\left(\frac{\omega_j}{\omega_i}\right)^{1/(2r)}
 \left(\frac{V_i}{V_j}\right)^{1/(2r)}.                         \tag{26}
\]

This is credited to PR \#25 and its source continuation. It uses the
original norm and volume sequence. It does not require constructing an
inverse before the first admissible degree.

Combining (1) and (26) proves the finite upper certificate

\[
 \boxed{
 \min_{i\le N<j}\epsilon_N
 \le\left(\frac{\omega_j}{\omega_i}\right)^{1/(2r)}
 \bigl((1+a)(1+b)^{q-1}\bigr)^{1/(2r)}.
 }                                                             \tag{27}
\]

For either quartet block \((i,j)=(q,2q)\) or \((q-1,2q-1)\), the
reviewed comparable-window norm bound supplies an inherited constant
\(C_h^{\mathrm{win}}\) independent of \(k\). Enlarge it once to account
for \(i=q-1\). Then

\[
 \boxed{
 \min_{i\le N<j}\epsilon_N
 \le C_h^{\mathrm{win}}q
 (1+a)^{1/(2q)}(1+b)^{(q-1)/(2q)}.
 }                                                             \tag{28}
\]

The selected member is still a canonical degree in the original window.
Its source relation, arithmetic unit, and action are not chosen again
after selecting the degree.

\subsubsection{Quantitative consequence of the already established
off-line lower
bound}\label{quantitative-consequence-of-the-already-established-off-line-lower-bound}

For the packet consisting exactly of an off-line quartet with full
common multiplicity, retain

\[
 q_k=[1+k(m-1)](k+1)^2,\qquad L_{h,k}\ge\delta kq_k/2.
\]

Let \(D_h>0\) be the fixed constant in the inherited per-block lower
inequality (2). The new upper bound implies, on each original block,

\[
 \boxed{(D_hk)^{2q}\le(1+a)(1+b)^{q-1}.}                        \tag{29}
\]

In the source area quantities it implies

\[
 \boxed{
 \frac{2\mathcal A_2}{(q-1)t_1}
 \ge\left(\frac{(D_hk)^{2q}}{1+t_1}\right)^{1/(q-1)}-1,
 }                                                             \tag{30}
\]

whenever that hypothetical lower budget is positive, so \(t_1>0\). When
the right side is negative the inequality remains valid but has no
positive force.

This strengthens the earlier statement that the mean relation cost must
grow: unless the total cost has an extremely large individual scale,
there must also be substantial \emph{two-direction exterior area} in the
original correction map. It is not a lower bound on every correction
eigenvalue.

For example, along such a family, a proved estimate
\(\log(1+t_1)=o(q\log k)\) would turn (30) into the necessary lower
exponent

\[
 \liminf\frac{\log(1+2\mathcal A_2/((q-1)t_1))}{\log k}\ge2.
\]

This last statement is a consequence with that displayed premise, not a
claimed arithmetic estimate. The actual unresolved calculation is to
bound the explicit sums (12)--(13) along the required growing-degree
quartet family. Their uniform size has not been established in this
note.

For the two-block total, (19) yields

\[
 \mathcal B_{h,k}\le U_2(Z^{(0)})+U_2(Z^{(1)}).
\]

The inherited lower threshold remains four. No new purity conclusion is
inferred from finite positivity or the optimality of a generic
determinant inequality.

\subsection{9. A finite certificate robust to moment
enclosures}\label{a-finite-certificate-robust-to-moment-enclosures}

The sharp formula (16) uses exact traces. One must not insert
uncertified decimal traces into it and call the result an upper
certificate.

There is a direct error-aware alternative. Suppose proved enclosures
give \(t_1\in[t_1^-,t_1^+]\) and \(t_2\in[t_2^-,t_2^+]\). A rational cap
\(A_0\ge t_1^+\) bounds every nonnegative eigenvalue. Choose a rational
\(0\le B_0<A_0\) and use the explicit polynomial \(P_{A_0,B_0}\) in
(20).

The same positive remainder proves

\[
 \log\det(I+Z)\le
 q\log(1+B_0)+\frac{t_1-qB_0}{1+B_0}
 +c_{A_0,B_0}(t_2-2B_0t_1+qB_0^2).                            \tag{31}
\]

Every coefficient in (31) is a rational combination of two logarithms of
positive rationals. Rational interval arithmetic, including all
products, gives a valid upper enclosure from the two trace intervals. It
needs no separately certified bound \(K_i\succeq g_0K_j\).

The checker encloses \(\log x\) by writing \(x=2^ey\), \(1\le y<2\), and
retaining both factors in

\[
 \log y=2\sum_{n=0}^{M-1}\frac{z^{2n+1}}{2n+1}+R_M,
 \quad z=\frac{y-1}{y+1},\quad
 0\le R_M\le\frac{2z^{2M+1}}{(2M+1)(1-z^2)}.
\]

The same formula encloses \(\log2\); reciprocation handles \(x<1\).
These finite interval checks certify their supplied rational data only.
An arithmetic application must supply certified moment/jet enclosures
and hence valid trace intervals.

\subsubsection{Arithmetic input size remains
explicit}\label{arithmetic-input-size-remains-explicit}

For a direct construction through degree \(j\), the source monic
coefficients and norms are determined by the original convolution
moments through degree \(2j\), and the quotient columns retain the full
coefficients of \(\chi\) and the specified arithmetic inclusion. The
sums (12)--(13) use \(G_i\); they do not require a second inverse
\(G_j\) or an independent spectral-gap proof. This is not a claim that
the required source moments are low-order or uniformly bounded. In the
two length-\(q\) endpoint blocks, this direct route can require moments
through \(4q\). The parallel session's separate lower-input upper-volume
comparison is not being replaced or credited as proved here.

\subsubsection{Coordinate and mass
propagation}\label{coordinate-and-mass-propagation}

Under a fixed invertible quotient-coordinate map \(C\), the actual data
become

\[
 F\mapsto CF,\quad G_i\mapsto C^{-*}G_iC^{-1},\quad
 Z\mapsto CZC^{-1}.
\]

Thus both traces, \(\mathcal A_2\), and the bound remain unchanged
through that similarity. Raw confluent derivative coordinates keep their
factorials in \(C\); the map is not silently made unitary.

A common source mass multiplier \(c>0\) gives

\[
 \Omega\mapsto c\Omega,\quad K_N\mapsto c^{-1}K_N,\quad
 G_N\mapsto cG_N,
\]

so \(Z,t_1,t_2,\mathcal A_2\) are unchanged by these displayed
composites. The original mass has not been reset. An ambient rectangular
arithmetic inclusion \(\eta\) is retained as an injection, not used as
an invertible square coordinate map.

\subsection{10. Exact source
calibrations}\label{exact-source-calibrations}

These are declared Gaussian polynomial fixtures, not asserted zeta
packets. Use original coordinate \(S=1+ix\), literal Gaussian mass
seven, variance one, and quotient \(\chi=(S-1)^q\). The monic source
recurrence and norms are

\[
 p_0=1,\ p_1=S-1,\quad p_{n+1}=(S-1)p_n+np_{n-1},\quad
 \omega_n=7n!.
\]

All quotient matrices below are computed in the original remainder basis
\(1,S,\ldots,S^{q-1}\).

For \(q=2\), the blocks \((1,3)\) and \((2,4)\) give

\[
 (t_1,t_2)=\left(2,\frac52\right),\qquad
 (t_1,t_2)=\left(\frac74,\frac{37}{16}\right).
\]

Their exact two-trace factors are \((a,b)=(3/2,1/2)\) and \((3/2,1/4)\),
so

\[
 \frac{V_1}{V_3}=\frac{15}{4},\quad
 \frac{V_2}{V_4}=\frac{25}{8},\quad
 \mathcal B=\log(375/32).
\]

The two-moment upper bound is exact in dimension two. This matches the
preceding restriction example without supplying a spectral gap.

For \(q=3\), the original block \((3,6)\) has

\[
 t_1=11,\quad t_2=\frac{3365}{32},\quad
 \mathcal A_2=\frac{507}{64},\quad \frac{V_3}{V_6}=\frac{5145}{256},
\]

\[
 a=\frac{11}{3}+\frac{7\sqrt{127}}{12},\qquad
 b=\frac{11}{3}-\frac{7\sqrt{127}}{24}.
\]

Numerical evaluations of these exact expressions are

\[
 \log(V_3/V_6)\approx3.000603203789,\quad
 U_2\approx3.063328308618,\quad
 3\log(1+t_1/3)\approx4.621335122841.
\]

The exact checks use rational enclosures and polynomial identities, not
these decimals as proof. The original operator has its full nonzero
nilpotent jets. No inference from a restriction eigenvalue to an
off-line arithmetic zero is made.

\subsection{11. What is established and what is handed
onward}\label{what-is-established-and-what-is-handed-onward}

Established here are the source cost map (7)--(10), the first and second
exterior source norms (12)--(15), the sharp two-trace upper bound with a
proved cap (16)--(21), its positive remainder and third-trace
refinement, the area-only upper certificate, and their composition with
the current first-degree-correct consecutive theorem. These are written
finite theorems; exact algebraic regression checks support their
implementation. They are not additional Lean certificates.

The same canonical arithmetic volume is now bounded without a separate
restriction-gap premise. Estimating the growing arithmetic values of
(12)--(13) remains a real analytic task. This note does not assert a
uniform bound completing RH, does not erase a non-radical packet, and
does not turn an \(e\)-valued quotient observation into absence.

No remote branch, publication cut, or formalization source was changed.
The handoff lists bounded reusable targets for the parallel session and
leaves its own interval-moment and fixed-budget work untouched.

\subsection{References and source
locators}\label{references-and-source-locators}

\begin{itemize}
\tightlist
\item
  Owner-supplied \emph{Tau Endpoint Restriction Control},
  \texttt{NOTE.tex} and original checker in the 13 September 2026 ZIP.
  Read as source; its finite map definitions are inputs here.
\item
  Owner-supplied \texttt{Pasted\ markdown(20260913-004825).md}: PR \#25
  report, norm-review clarification, exact product, first-degree step,
  and threshold-four result. Source claims are distinguished from this
  session's executions.
\item
  Owner-supplied \texttt{Pasted\ markdown\ (2)(5).md}: ongoing
  collaboration transcript; partial reported results are not substituted
  for their missing full proofs.
\item
  GitHub \texttt{KokunoYumeto/zeta-function-research-reader}, main
  \texttt{952ef9fee1e1419b6858d920354de8fa99430b7d}; PR \#25
  \texttt{66fac5e7885a40cd4873e74b754907320f05b067},
  \texttt{workbenches/tau-consecutive-window-formal/RESEARCH\_NOTE.md}.
\item
  Grone, B., Johnson, C. R., Marques de Sá, E., and Wolkowicz, H.
  (1984). \emph{Improving Hadamard's inequality}. Linear and Multilinear
  Algebra 16(1--4), 305--322. DOI: 10.1080/03081088408817634. The
  publisher abstract and author-hosted bibliography were inspected; the
  present proof is given explicitly rather than attributed to an unread
  theorem number.
\item
  NIST Digital Library of Mathematical Functions, §3.3, especially the
  divided-difference and confluent interpolation formulas.
  https://dlmf.nist.gov/3.3
\end{itemize}

\end{document}
